
\section{Background}
\label{sec:background}

In ASCA, keystroke classification relies on transforming raw audio signals into representations like mel-spectrograms, and training models to identify the acoustic patterns within them. Recent developments in deep learning offer architectures specifically tailored for image data, including spectrograms~\cite{yang2023slnet, esmaeilpour2020sound, taheritajar2024survey}. Several key architectures are discussed in this paper and introduced in the following.

\textbf{CoAtNet} (Convolutional Attention Network)~\cite{dai2021coatnet} is a hybrid architecture that integrates Convolutional Neural Networks (CNNs) with self-attention mechanisms. It leverages the strength of both architectures. The convolutional layers are suitable for extracting local spatial features, e.g., edges and textures in the mel-spectrogram for keystroke audios. The self-attention mechanism captures long-range dependencies across the entire input. It focuses on the temporal relationships in the keystrokes. CoAtNet's capability of analyzing local and long-term features is promising for the ASCA task. Its staged design by stacking convolutional blocks before transformer layers enables efficient hierarchical feature extraction.

\textbf{Swin Transformer} enhances Vision Transformers (VTs) by introducing a hierarchical, shifted-window approach to self-attention~\cite{liu2021swin}. Unlike standard VTs that treat images as sequences of patches and apply global self-attention, Swin partitions images into non-overlapping patches and computes self-attention within local windows, significantly reducing computational complexity. It periodically shifts these windows, enabling partial overlaps that facilitate global context aggregation. Furthermore, Swin Transformer adjusts patch and window sizes at different stages, allowing it to efficiently capture varying levels of detail and improving both scalability and performance in vision tasks.

\textbf{Audio models,} particularly those designed for speech recognition, such as OpenAI Whisper~\cite{radford2023robust} and Mozilla DeepSpeech~\cite{hannun2014deep}, are primarily built for tasks like speech-to-text or text-to-speech conversion. While these models excel at processing linguistic content, they are not optimized for detecting short-duration, transient events like keystrokes. As a result, their pre-trained models cannot be directly applied to the task of ASCA on keyboards unless a large dataset containing long sequences of keystrokes is available.

\textbf{Architecturally, these models share similarities with the method proposed in this work}, as they also transform audio signals into mel-spectrograms and analyze the resulting visual representations. However, due to the need to process a significantly larger vocabulary, these models are inherently heavyweight. Consequently, this work does not utilize general-purpose audio-to-text models.